\documentclass[11pt,a4paper]{article}
% Artículo de resultados. Las cifras vienen de ../presentacion/resultados.tex,
% que genera scripts/09_presentacion.py a partir de data/processed/ (igual que la presentación).
\usepackage[utf8]{inputenc}
\usepackage[T1]{fontenc}
\usepackage[spanish,es-nodecimaldot,es-noshorthands,es-tabla]{babel}
\usepackage{lmodern}
\usepackage[a4paper,margin=2.4cm]{geometry}
\usepackage{microtype}
\usepackage{booktabs}
\usepackage{graphicx}
\usepackage{xcolor}
\usepackage{amsmath}
\usepackage[font=small,labelfont=bf]{caption}
\usepackage[section]{placeins} % las figuras no pasan a la sección siguiente
\usepackage[hidelinks]{hyperref}
\graphicspath{{../figures/}}

\newcommand{\pendiente}{\textcolor{orange}{[pendiente]}}
% Archivo generado por scripts/09_presentacion.py: no editar a mano
\newcommand{\periodoMeteo}{2025-01-01 a 2026-10-31}
\newcommand{\umbralParacas}{10}
\newcommand{\ventanaParacas}{12--18}
\newcommand{\mesCandidato}{2025-01}
\newcommand{\eventosCandidato}{10}
\newcommand{\criterioCandidato}{fijado en config.yaml (vientos\_paracas.mes\_candidato)}
\newcommand{\fuenteCandidato}{METAR}
\newcommand{\eventosMetar}{101}
\newcommand{\eventosEra}{0}
\newcommand{\mesMaxMetar}{2025-01 (10 días)}
\newcommand{\vientoFuerteObs}{11.0}
\newcommand{\vientoFuerteEra}{6.4}
\newcommand{\tablaValidacion}{velocidad (m/s) & 15984 & 0.20 & 1.96 & 0.59 \\
velocidad tarde 12-18 h (m/s) & 4665 & -0.46 & 2.29 & 0.36 \\
velocidad cuando obs > 10 m/s & 287 & -4.57 & 4.68 & 0.48 \\
dirección (grados) & 12399 & -25.28 & 52.61 & -- \\
temperatura (°C) & 15981 & -0.25 & 1.18 & 0.92 \\}
\newcommand{\seriePrincipal}{UHSLC}
\newcommand{\anioUhslc}{2022}
\newcommand{\qcAprobados}{1 de 16}
\newcommand{\pendienteIoc}{0.73}
\newcommand{\factorF}{0.75}
\newcommand{\claseMarea}{mixta, predominantemente semidiurna}
\newcommand{\ampMdos}{31.1}
\newcommand{\ampSdos}{9.2}
\newcommand{\ampNdos}{8.8}
\newcommand{\ampKuno}{20.0}
\newcommand{\ampOuno}{10.3}
\newcommand{\ampPuno}{6.1}
\newcommand{\rangoSicigia}{0.81}
\newcommand{\rangoCuadratura}{0.44}
\newcommand{\rangoMaximo}{1.41}
\newcommand{\rangoSicigiaObs}{1.32}
\newcommand{\rangoCuadraturaObs}{0.72}
\newcommand{\varExplicada}{93.8}
\newcommand{\rmsResiduo}{7.6}
\newcommand{\seriePrincipalNombre}{UHSLC 2022}
\newcommand{\tablaArmonicos}{M2 & 12.42 & 31.0 $\pm$ 0.1 & 337.5 $\pm$ 0.2 & 382060 \\
K1 & 23.93 & 20.0 $\pm$ 0.1 & 28.2 $\pm$ 0.3 & 94660 \\
O1 & 25.82 & 10.3 $\pm$ 0.1 & 347.3 $\pm$ 0.7 & 22187 \\
S2 & 12.00 & 9.2 $\pm$ 0.1 & 350.2 $\pm$ 0.5 & 35827 \\
N2 & 12.66 & 8.8 $\pm$ 0.1 & 297.8 $\pm$ 0.7 & 21084 \\
P1 & 24.07 & 6.1 $\pm$ 0.1 & 26.4 $\pm$ 1.0 & 10813 \\}
\newcommand{\dAmMdosIoc}{-0.9}
\newcommand{\dgMdosIoc}{4.3}
\newcommand{\areaBahia}{\pendiente}
\newcommand{\volBahia}{\pendiente}
\newcommand{\profMedia}{\pendiente}
\newcommand{\profMax}{\pendiente}
\newcommand{\dAmMdosFes}{\pendiente}
\newcommand{\dgMdosFes}{\pendiente}


% Figura con aviso si aún no existe el PNG
\newcommand{\figura}[2]{%
  \IfFileExists{../figures/#1.png}%
    {\includegraphics[#2]{#1}}%
    {\fbox{\parbox[c][3cm][c]{0.8\linewidth}{\centering\small\textcolor{orange}{Figura pendiente:}\\\texttt{\detokenize{figures/#1.png}}}}}}

\title{\textbf{Viento Paracas, marea y circulación forzada por el viento en la bahía de Paracas (Perú)}\\[0.4em]\large Caracterización con observaciones, reanálisis y AEM3D}
\author{Curso de Oceanografía\\Universidad Nacional Mayor de San Marcos}
\date{\today}

\begin{document}
\maketitle

\begin{abstract}
\noindent
Se caracterizan el viento, la marea y la batimetría de la bahía de Paracas y se simula su respuesta a la
marea y al viento con el modelo hidrodinámico AEM3D. Entre \periodoMeteo, la estación METAR del aeropuerto
de Pisco (SPSO) registra \eventosMetar{} días con viento Paracas (más de \umbralParacas~m/s entre las
\ventanaParacas~h locales), mientras que ERA5 no registra ninguno: el reanálisis reproduce la dirección y el
ciclo diurno, pero da \vientoFuerteEra~m/s cuando se observan \vientoFuerteObs~m/s. Una corrección por
cuantiles iguala el número de eventos, pero no su cronología, por lo que el forzamiento usa el viento observado.
La marea es \claseMarea{} ($F$ = \factorF), con un rango de sicigia observado de \rangoSicigiaObs~m. El registro
del mareógrafo de Pisco tiene un error de escala ($\times$\razonError) entre mayo de 2016 y enero de 2025, que se
detectó comparándolo con Callao, San Juan y FES2022b; corregido ese error, FES2022b reproduce la marea local.
Una simulación de 14 días (\aemInicio{} al \aemFin) con agua homogénea reproduce la marea del mareógrafo
(RMSE \aemRmseIoc~cm) y muestra que en las tardes con evento Paracas la corriente de superficie de la bahía
sube de \corrCalma{} a \corrEvento~cm/s, con un retraso de \rezagoVientoCorr~h respecto del viento
($r$ = \rVientoCorr). La corriente promediada en la vertical es unas cinco veces menor: el viento mueve una
capa delgada de superficie.

\medskip\noindent\textbf{Palabras clave:} bahía de Paracas, viento Paracas, marea, FES2022, ERA5, AEM3D.
\end{abstract}

% ===================================================================
\section{Introducción}

La bahía de Paracas, en la costa central del Perú, es una bahía somera, abierta al norte y protegida al oeste
por la península de Paracas. Su clima está marcado por los vientos Paracas: vientos fuertes del sur y suroeste
que se intensifican por la tarde y que pueden levantar arena y oleaje dentro de la bahía. Para estudiar cómo
circula el agua de la bahía hacen falta tres cosas: un forzamiento atmosférico que represente bien esos
eventos, una marea confiable en las fronteras abiertas y una batimetría razonable.

Este trabajo arma esos tres componentes con datos públicos y verificados, evalúa sus errores y los usa para
forzar una simulación con AEM3D \cite{aem3d}. Las preguntas son: (1) ¿representa ERA5 los vientos Paracas?;
(2) ¿qué tan confiables son el mareógrafo de Pisco y el modelo global de marea FES2022b?; y (3) ¿cuánto cambia
la corriente de la bahía durante un evento Paracas?

% ===================================================================
\section{Área de estudio y datos}

El dominio cubre 13.60°--13.95°S y 76.15°--76.50°W (figura~\ref{fig:area}). Se usaron:

\begin{itemize}
  \item \textbf{ERA5} \cite{era5}, horario, en el punto más cercano a la boca de la bahía: viento a 10~m,
    temperatura y punto de rocío a 2~m, presión, radiación de onda corta y larga, precipitación y nubosidad.
  \item \textbf{METAR SPSO} (aeropuerto de Pisco, 13.745°S, 76.220°W), del Iowa Environmental Mesonet
    \cite{iem}, convertido a unidades SI y con control de calidad (duplicados, valores imposibles, calmas).
  \item \textbf{Mareógrafo de Pisco} (13.819°S, 76.252°W): datos en tiempo real del IOC \cite{ioc} y
    datos horarios de calidad de investigación del UHSLC (estación 683) \cite{uhslc}.
  \item \textbf{Mareógrafos vecinos} de Callao (UHSLC 093) y San Juan de Marcona (UHSLC 096), como control.
  \item \textbf{FES2022b} \cite{fes2022}, solo elevaciones de marea oceánica, evaluado con pyTMD \cite{pytmd}.
  \item \textbf{GEBCO\_2026} \cite{gebco}, batimetría en una malla de 15 segundos de arco.
\end{itemize}

\begin{figure}[t]
  \centering\figura{batimetria_mapa}{width=0.75\textwidth}
  \caption{Área de estudio: batimetría de GEBCO\_2026, polígono de la bahía, mareógrafo de Pisco y estación SPSO.}
  \label{fig:area}
\end{figure}

% ===================================================================
\section{Métodos}

\subsection{Viento}
Se define un \emph{evento Paracas} como un día con viento mayor de \umbralParacas~m/s entre las
\ventanaParacas~h locales (UTC$-$5). ERA5 se comparó con SPSO hora a hora (sesgo, RMSE y correlación; para
la dirección, la diferencia circular). Para corregir ERA5 se probó un ajuste por cuantiles
\cite{panofsky,maraun}, evaluado con validación cruzada que deja fuera un mes cada vez. El mes candidato
para modelar es \mesCandidato, con \eventosCandidato{} eventos según SPSO; empata con marzo de 2025, pero en
marzo no hay datos del mareógrafo IOC para validar la marea.

\subsection{Marea}
Los datos del IOC se limpiaron de picos, saltos y valores planos y se promediaron a valores horarios. El
análisis armónico se hizo con UTide \cite{utide,ttide} sobre un año completo. Para detectar errores de
escala se comparó, año a año, la amplitud de M$_2$ de Pisco con la de FES2022b y con la de los mareógrafos de
Callao y San Juan, y el nivel medio mensual de Pisco con el de Callao. El residuo se analizó con el método de
Welch \cite{welch}.

\subsection{Batimetría}
GEBCO se recortó al dominio. Dentro del polígono de la bahía se calcularon área, volumen y profundidad media.

\subsection{Modelo AEM3D}
Se armó una malla de \mallaDx~m (\mallaFilas{} $\times$ \mallaColumnas{} celdas, \mallaAgua{} columnas de agua)
con \mallaCapas{} capas estiradas: 1~m de espesor hasta 25~m de profundidad y más gruesas hasta
\mallaProfMax~m (figura~\ref{fig:malla}). Las fronteras norte, oeste y sur son abiertas y reciben la marea de
FES2022b. La superficie recibe el viento observado de SPSO y la temperatura, humedad, radiación y presión de ERA5 en la
boca de la bahía. El agua es
homogénea: sin estratificación y sin intercambio de calor con la atmósfera. En esta versión de AEM3D el
viento solo transfiere momento al agua a través del modelo de capa de mezcla, por lo que se activó la
densidad y la temperatura con temperatura uniforme. La versión de demostración limita la simulación a 14 días:
se simuló del \aemInicio{} al \aemFin{} (hora local), con paso de 30~s, y se descartaron las 3 primeras horas.

\begin{figure}[t]
  \centering\figura{aem3d_malla}{width=0.85\textwidth}
  \caption{Malla de AEM3D: profundidad de cada celda y espesor de las capas.}
  \label{fig:malla}
\end{figure}

% ===================================================================
\section{Resultados}

\subsection{Viento: ERA5 subestima los eventos Paracas}

El viento es del S--SO, con un fuerte ciclo diurno: máximo por la tarde y mínimo de madrugada. En promedio, ERA5
se parece a SPSO (tabla~\ref{tab:validacion}), pero no reproduce los extremos: cuando SPSO pasa de 10~m/s, ERA5
da en promedio \vientoFuerteEra~m/s. En todo el periodo, SPSO registra \eventosMetar{} días con evento y
ERA5, \eventosEra{} (figura~\ref{fig:eventos}).

\begin{table}[h]
  \centering\small
  \caption{ERA5 frente a SPSO, \periodoMeteo. Sesgo = ERA5 $-$ observado.}
  \label{tab:validacion}
  \begin{tabular}{@{}lrrrr@{}}
    \toprule
    Variable & $n$ & Sesgo & RMSE & $r$ \\
    \midrule
    \tablaValidacion
    \bottomrule
  \end{tabular}
\end{table}

La corrección por cuantiles iguala el total de eventos (\eventosQmTotal{} frente a \eventosMetar), pero los
coloca en los meses equivocados: la correlación entre los conteos mensuales es \corrEventosQm{}, y en
\mesCandidato{} da \eventosQmMes{} evento en lugar de \eventosCandidato. Por eso el forzamiento usa el viento
observado de SPSO (\horasSpso{} de \horasForzamiento~h; las \horasRellenas~h faltantes se rellenan con ERA5
corregido).

\begin{figure}[t]
  \centering\figura{validacion_eventos_paracas_mensual}{width=\textwidth}
  \caption{Días con evento Paracas por mes según SPSO y ERA5.}
  \label{fig:eventos}
\end{figure}

\subsection{Marea: un error de escala en el mareógrafo de Pisco}

La M$_2$ de Pisco medida entre \aniosError{} es \razonError{} veces mayor que la de \aniosBuenos{}, mientras que
en Callao y San Juan no cambia. FES2022b coincide con la M$_2$ de Callao y San Juan y con la del Pisco de \aniosBuenos. En 2022,
la razón entre la amplitud de FES2022b y la observada es \razonFesCallao{} en Callao y \razonFesSanJuan{} en
San Juan, pero \razonFesPisco{} en Pisco, y lo mismo ocurre en todas las constituyentes
(figura~\ref{fig:vecinos}). El salto aparece justo después de un hueco de abril de 2016 y sigue
en enero de 2025, de modo que se trata de un error de escala del registro y no de un cambio de la marea. Por eso
la serie principal es el IOC de \ventanaPrincipal{} (\completitudPrincipal~\% de horas válidas), con UHSLC
\anioUhslc{} como referencia; las dos coinciden en M$_2$ (\dAmMdosIoc~cm y \dgMdosIoc°).

\begin{figure}[t]
  \centering\figura{marea_control_vecinos}{width=\textwidth}
  \caption{Control con mareógrafos vecinos. Izquierda: amplitud anual de M$_2$ en Callao, Pisco y San Juan (UHSLC), con la de FES2022 en cada estación (líneas discontinuas) y la del IOC 2025--26 en Pisco (rombo); en rosa, los años de Pisco que se apartan más de un 10~\% de sus vecinos. Derecha: amplitud de FES2022 dividida por la observada en 2022.}
  \label{fig:vecinos}
\end{figure}

El nivel medio mensual de Pisco sube unos \rangoNivelMedio~cm entre fines de 2025 y mediados de 2026. Callao
muestra la misma subida ($r$ = \corrNivelMedio, diferencia máxima \difNivelMedio~cm), así que es una señal
oceanográfica regional y no una deriva del sensor (figura~\ref{fig:nivelmedio}).

\begin{figure}[t]
  \centering\figura{marea_nivel_medio_pisco_callao}{width=0.85\textwidth}
  \caption{Nivel medio mensual en Pisco (IOC) y Callao.}
  \label{fig:nivelmedio}
\end{figure}

La marea es \claseMarea{} ($F = (K_1+O_1)/(M_2+S_2)$ = \factorF{} \cite{pugh}; tabla~\ref{tab:armonicos}). Los rangos
observados de sicigia y cuadratura son \rangoSicigiaObs{} y \rangoCuadraturaObs~m, mayores que los que dan
solo M$_2$ y S$_2$ (\rangoSicigia{} y \rangoCuadratura~m) por el aporte diurno; el rango máximo es
\rangoMaximo~m. Los armónicos explican el \varExplicada~\% de la varianza y el residuo tiene un RMS de
\rmsResiduo~cm. FES2022b reproduce la M$_2$ del mareógrafo ($\Delta A$ = \dAmMdosFes~cm,
$\Delta g$ = \dgMdosFes°) y, en \mesCandidato, el RMSE frente al IOC corregido baja de \rmseCandSin{} a
\rmseCandCorr~cm al aplicar el factor \factorCandidato.

\begin{table}[h]
  \centering\small
  \caption{Constituyentes principales, IOC \ventanaPrincipal. Intervalos del 95~\%; fase de Greenwich (UTC).}
  \label{tab:armonicos}
  \begin{tabular}{@{}lrrrr@{}}
    \toprule
    Const. & Periodo (h) & $A$ (cm) & $g$ (°) & SNR \\
    \midrule
    \tablaArmonicos
    \bottomrule
  \end{tabular}
\end{table}

\subsection{Batimetría}
Dentro del polígono, la bahía tiene \areaBahia~km$^2$, \volBahia~hm$^3$, una profundidad media de
\profMedia~m y una máxima de \profMax~m.

\subsection{Simulación con AEM3D}

\paragraph{Marea.} AEM3D transmite la marea de las fronteras al interior de la bahía sin deformarla: en la celda
del mareógrafo, el RMSE frente a FES2022b es \aemRmseFes~cm. Frente al IOC observado y corregido, el RMSE es
\aemRmseIoc~cm ($r$ = \aemRIoc, razón de amplitudes \aemRazonIoc), igual que el de FES2022b solo
(\fesRmseIoc~cm; figura~\ref{fig:aemmarea}).

\begin{figure}[t]
  \centering\figura{aem3d_marea_mareografo}{width=\textwidth}
  \caption{Nivel del mar en el mareógrafo de Pisco: IOC observado y corregido, FES2022b y AEM3D.}
  \label{fig:aemmarea}
\end{figure}

\paragraph{Corrientes.} La tabla~\ref{tab:corrientes} compara las tardes (12 a 18~h) con y sin evento Paracas.
En las tardes con evento, la corriente media de superficie es un 50~\% mayor y se dirige hacia el NE
(\dirEvento°), en el sentido general del viento del S--SO (figura~\ref{fig:corrientes}). La rapidez de la corriente sigue al
viento con \rezagoVientoCorr~h de retraso ($r$ = \rVientoCorr; figura~\ref{fig:vientocorriente}). Promediada
en la vertical, la corriente es unas cinco veces menor que en la superficie. La velocidad máxima de superficie
en todo el dominio es de \aemVelMax~cm/s, en mar abierto.

\begin{table}[h]
  \centering\small
  \caption{Corrientes medias dentro del polígono de la bahía, tardes de 12 a 18~h.}
  \label{tab:corrientes}
  \begin{tabular}{@{}lrrr@{}}
    \toprule
    & Superficie, media & Superficie, rapidez instantánea & Promedio vertical \\
    & (cm/s) & (cm/s) & (cm/s) \\
    \midrule
    Con evento Paracas & \corrEvento & \rapEvento & \corrVertEvento \\
    Sin evento & \corrCalma & \rapCalma & \corrVertCalma \\
    \bottomrule
  \end{tabular}
\end{table}

\begin{figure}[t]
  \centering\figura{aem3d_corrientes_eventos}{width=\textwidth}
  \caption{Corriente media de superficie simulada en las tardes con y sin evento Paracas. Línea roja: polígono de la bahía.}
  \label{fig:corrientes}
\end{figure}

\begin{figure}[t]
  \centering\figura{aem3d_viento_corriente}{width=\textwidth}
  \caption{Viento de SPSO (forzamiento) y rapidez media de la corriente de superficie en la bahía. En rosa, las tardes con evento Paracas.}
  \label{fig:vientocorriente}
\end{figure}

% ===================================================================
\section{Discusión}

\paragraph{El forzamiento atmosférico.} La subestimación de ERA5 es consistente con su resolución
(unos 30~km): el reanálisis no resuelve la canalización local del viento entre la península y el continente.
La corrección por cuantiles arregla la distribución de velocidades, pero no puede crear eventos en los días en
que ERA5 no tiene la señal. Para simular eventos concretos hace falta, por ahora, el viento observado. SPSO
está en tierra, así que sobre el agua el viento podría ser algo mayor y las corrientes simuladas serían una
estimación conservadora.

\paragraph{La calidad del registro de marea.} El error de escala de Pisco no se nota mirando solo la serie:
apareció al comparar con FES2022b y con mareógrafos vecinos. Sin ese control, el análisis armónico habría dado
amplitudes un 35~\% mayores y se habría concluido, erróneamente, que FES2022b subestima la marea local.

\paragraph{La respuesta de la bahía.} Con agua homogénea, el viento acelera sobre todo una capa delgada de
superficie, con flujo de retorno debajo. Con una termoclina, el momento del viento quedaría atrapado en una capa
de mezcla aún más delgada, por lo que la corriente de superficie probablemente sería mayor. La corrida baroclínica
necesita perfiles iniciales de temperatura y salinidad, y perfiles en las fronteras, que todavía no se tienen.

\paragraph{Limitaciones.} La versión de demostración de AEM3D limita la simulación a 14 días, por lo que quedan
fuera los eventos del 7 al 9 de enero de 2025. GEBCO es grueso dentro de la bahía; la fusión con los sondajes de
la carta náutica de la DHN está preparada pero sin datos.

% ===================================================================
\section{Conclusiones}

\begin{enumerate}
  \item El viento de la bahía es del S--SO, con fuerte ciclo diurno; los vientos Paracas ocurren por la tarde.
    ERA5 reproduce la dirección y el ciclo, pero no la intensidad de los eventos (\eventosEra{} frente a
    \eventosMetar{} días), y la corrección por cuantiles no recupera su cronología.
  \item La marea de Pisco es \claseMarea, con M$_2$ = \ampMdos~cm y K$_1$ = \ampKuno~cm. El registro tiene un
    error de escala de $\times$\razonError{} entre 2016 y enero de 2025; corregido, FES2022b reproduce la marea
    local y sirve como forzamiento de frontera.
  \item En AEM3D, la corriente de superficie de la bahía sube de \corrCalma{} a \corrEvento~cm/s en las tardes
    con evento Paracas y responde al viento en \rezagoVientoCorr~h. El viento actúa sobre una capa delgada de
    superficie.
\end{enumerate}

\section*{Disponibilidad del código y los datos}
Todos los scripts, la configuración y la presentación están en el repositorio del proyecto. Los datos se
descargan o regeneran con los scripts (ver \texttt{README.md}).

% ===================================================================
\begin{thebibliography}{99}
\small
\bibitem{era5} Hersbach, H., et al. (2020). The ERA5 global reanalysis. \emph{Q. J. R. Meteorol. Soc.}, 146, 1999--2049. doi:10.1002/qj.3803
\bibitem{iem} Iowa Environmental Mesonet. ASOS/METAR, estación SPSO. Iowa State University.
\bibitem{ioc} VLIZ; IOC. Sea level station monitoring facility. doi:10.14284/482
\bibitem{uhslc} Caldwell, P. C., Merrifield, M. A., Thompson, P. R. (2015). Joint Archive for Sea Level holdings (UHSLC). NOAA NCEI. doi:10.7289/V5V40S7W
\bibitem{fes2022} AVISO. FES2022b, elevaciones de marea oceánica. doi:10.24400/527896/a01-2024.004
\bibitem{gebco} GEBCO Compilation Group (2026). GEBCO\_2026 Grid. doi:10.5285/4f68d5c7-45eb-f999-e063-7086abc036fa
\bibitem{utide} Codiga, D. L. (2011). Unified tidal analysis and prediction using the UTide Matlab functions. GSO Tech. Rep. 2011-01, URI.
\bibitem{ttide} Pawlowicz, R., Beardsley, B., Lentz, S. (2002). Classical tidal harmonic analysis including error estimates in MATLAB using T\_TIDE. \emph{Comput. Geosci.}, 28, 929--937.
\bibitem{pytmd} Sutterley, T. C., et al. pyTMD: Python-based tidal prediction software (v3.1).
\bibitem{pugh} Pugh, D., Woodworth, P. (2014). \emph{Sea-Level Science}. Cambridge University Press.
\bibitem{welch} Welch, P. D. (1967). The use of fast Fourier transform for the estimation of power spectra. \emph{IEEE Trans. Audio Electroacoust.}, 15(2), 70--73.
\bibitem{panofsky} Panofsky, H. A., Brier, G. W. (1968). \emph{Some Applications of Statistics to Meteorology}. Penn State University.
\bibitem{maraun} Maraun, D. (2016). Bias correcting climate change simulations. \emph{Curr. Clim. Change Rep.}, 2, 211--220.
\bibitem{aem3d} Hodges, B., Dallimore, C. \emph{AEM3D User Manual} y \emph{Science Manual}. Hydronumerics.
\end{thebibliography}

\end{document}
